\section{Conclusion}
\label{sec:conclusion}

VBVR-Pro was built to ask whether generative models can reason in the visual medium. We kept its tasks, its inputs, its output container and its evaluator, and changed only how the answer is produced: a coding agent writes a program that renders the video instead of denoising it. On VBVR-Pro-Bench that one change moves the best score from 0.670, for a video model RL-trained on the benchmark's own in-domain families, to 0.923; all three closed-model agents exceed 0.89. The direction of generalisation flips as well. Trained video models fall on the held-out families and coding agents rise on them, so the agents' advantage is largest exactly where the video models were not trained. Among open-weight models the best reaches 0.572, and the weakest fail before reasoning begins, by never using a tool or by giving up after one or two calls. The lesson is not that video models cannot reason, but that on tasks whose answers are programs, a solver that writes the program is scored on the reasoning alone, while a solver that paints the pixels is scored on the reasoning and the painting at once.

\paragraph{Limitations.}
The official evaluator is used unmodified and inherits its properties. On a fresh seed a frozen first frame scores above 0.5 on four of the 100 task families and the generator's own ground truth scores below 0.95 on six, so a small part of every score is evaluator noise. All results are on one benchmark whose tasks are generated and rendered by code, the setting in which writing code is most natural. The comparison is between an agent with tools, a sandbox and up to 30 minutes of iteration and a video model that emits one sample in one shot; it measures what each route delivers, not what each model could do under the other's constraints. Agent cost is reported in tokens and wall-clock, not dollars, and varies by two orders of magnitude across models. The natural next step is to train the code-writing policy with GRPO \citep{shao2024deepseekmath} using the official evaluator as the reward, mirroring the protocol VBVR-Pro applied to video models \citep{vbvrpro2026}; the reward pipeline and a training step on AWS Trainium are in place, and no result from it is reported here.
